\section{Data and code availability}
\label{sec:availability}

The implementation repository contains the development controller, the protected assurance path, the local hostile falsifier suite, the hidden-cause case suite, the diagnostic attribution archive and the table generator. Every artifact named in the narrative is listed here with its path, so that no repository location has to be carried in the prose. Source, protocols and evidence are version-controlled in the \idt{SzeChunYiu/ORION} repository, under \idt{papers/orion-15-self-orion/}; persistent DOI assignment remains pending publication.

\subsection*{Preregistration}

The original protected-campaign design was frozen before that campaign accessed
any outcome, under the identifier
\texttt{P5.hidden-cause-fresh-transfer.v1}. Its recorded state is a design
freeze, not an execution freeze. A reader checking the prospective margins in
Section~\ref{sec:methods} should compare them against that artifact. This timing
claim does not extend to every later successor: P5-RD-02 has separate V1/V2
development identities, and V3 was frozen after V2 outcomes were known for an
audit-only archival replay.

\subsection*{Frozen design, case structure and acceptance rules}

The design freeze is \idt{protocol/PROTOCOL_V1.json}; the hidden-cause case structure is \idt{protocol/HIDDEN_CAUSE_CASE_SCHEMA_V1.json}; replay/freshness definitions and harm accounting are \idt{protocol/FRESH_TRANSFER_POLICY_V1.md}; and the staged acceptance rules and their successor revision are \idt{protocol/PROTOCOL_V2.json} and \idt{protocol/STAGED_ACCEPTANCE_POLICY_V2.md}. The protected-suite freeze procedure is \idt{protocol/PROTECTED_SUITE_FREEZE_V1.md}, and the retention chain for negative history is \idt{protocol/P5_NEGATIVE_HISTORY_CHAIN_V1.json}. The revision-level successor benchmark of Section~\ref{sec:methods} is frozen as \idt{protocol/SELF_ORION_V3_REVISION_LEVEL_PROTOCOL_V1.json}; the additive metadata-only correction that binds MDA v3 structurally while retaining the unreleased-solver blocker is \idt{protocol/SELF_ORION_V3_MDA_SOURCE_BINDING_AMENDMENT_2026-08-23.json}. The external transfer benchmark used as a reference point is pinned to \idt{Gen-Verse/PAST-Bench@f8223517ae7491e776b69793d9f11e9d074ab42e}; any use of it still requires a frozen task and split manifest.

The additive wide design is
\idt{protocol/P5_WIDE_REVISION_LEVEL_SUCCESSOR_V1.json}; its three
error-derived identification studies are
\idt{protocol/P5_RESIDUAL_DISCRIMINATOR_SUCCESSORS_V1.json}. Those protocol
files record their outcome-blind state at their own design-origin timestamps;
they are not a current claim that every descendant study remains unexecuted.
The wide campaign, P5-RD-01 and P5-RD-03 remain unexecuted. P5-RD-02's later
public-development V1/V2 observations and post-outcome V3 archival audit are
listed below. None modifies the original preregistration, corrects a diagnostic
error, or turns public replay into protected confirmatory evidence.
The provider- and modality-diverse candidate frame is
\idt{development/provider-diverse-metadata-census-2026-08-23/SOURCE_CENSUS_V1.json},
with protocol
\idt{development/provider-diverse-metadata-census-2026-08-23/CENSUS_PROTOCOL_V1.json}
and receipt
\idt{development/provider-diverse-metadata-census-2026-08-23/VERIFICATION_RECEIPT_V1.json}.
These artifacts contain no project/issue cases, labels, outputs or outcomes;
their rights and eligibility terminals remain fail-closed.
The outcome-blind six-arm identity and public-panel preflight is archived under
\idt{development/p5-six-arm-comparator-panel-2026-08-23/}. Its protocol,
normalized identity ledger, public-source rights matrix, retained
negative-result ledger, primary-source research sidecars and checksum manifest
distinguish semantic/source identity from execution readiness and record each
official-versus-reimplementation status, licence byte, entrypoint,
information/action/cost envelope, native error terminal and missing adapter.
No comparator performance, dataset row or protected outcome is present.
The terminal-preserving adapter successor is archived under
\idt{development/p5-six-arm-terminal-adapters-v2-2026-08-23/}. Its prospective
protocol and six same-visible-symptom fixtures were frozen before any native
output example; the per-arm ledgers, resource/configuration audit,
factorization theorem, deterministic conformance program, 90-case receipt,
recursive negative-result ledger and \idt{SHA256SUMS} retain native failures
and unsupported or mixed fibres as \texttt{UNRESOLVED}. The archive contains no
comparator execution, native output example, public or protected outcome, or
performance estimate. It leaves zero of six arms confirmatory-ready and
preserves H1--H4 and protected freshness as \textsf{CANNOT\_CHECK}.
The stricter outcome-free refinement is archived under
\idt{development/p5-six-arm-adapter-refinement-v3-2026-08-23/}. Its frozen
certificate schema, action/write-surface schema, eight-class front registry,
native-terminal rules and deterministic validator cover exactly 231 declared
synthetic cases. The receipt records 40 typed case records in 20 supported
constant fibres and 191 \texttt{UNRESOLVED}, with zero raw-native licences and
zero conformance failures. The separate execution ledger preserves 108 exact
blocker instances, 18 per arm, and zero of six execution-ready. Its
\idt{SHA256SUMS} also pins the unchanged V2 predecessor manifest. This is
contract conformance only: no comparator, native example, public/protected
outcome or performance table was opened.

The outcome-blind C1 SWE-agent execution-binding successor is archived under
\idt{development/p5-swe-agent-execution-binding-v4-2026-08-23/}. Its field
registry, pinned source/rights manifest, native parser and output schema,
fail-closed write-surface wrapper, resource/dependency lock, candidate-case and
custody schemas, adverse-result ledger, audit receipt and \idt{SHA256SUMS}
bind 9/21 C1 fields and preserve 12 blockers. The exact terminal is
\idt{P5_C1_V4_SWE_AGENT_NATIVE_PARSER_AND_ISOLATION_POLICY_BOUND__TWELVE_C1_FIELDS_BLOCKING__ZERO_OF_SIX_PANEL_READY__PERFORMANCE_AND_SUPERIORITY_CANNOT_CHECK}.
The packet copies only an outcome-free 101-byte upstream smoke fixture and no
substantive case of its own. No comparator/model, protected scorer, gold patch,
protected-test value, performance outcome or superiority result was opened.

The distinct outcome-blind C2 MOSS execution-binding successor is archived
under \idt{development/p5-moss-execution-binding-v4-2026-08-23/}. Its field
registry, exact source/rights manifest, native parser and terminal rules,
fail-closed runner, host and OpenClaw locks, smoke receipt, adverse-result
ledger and 70-check audit bind 7/21 C2 fields and preserve 14 blockers. The
exact terminal is
\idt{P5_C2_V4_MOSS_NATIVE_PARSER_DEPENDENCY_LOCKS_FALLBACK_AND_WALLCLOCK_BOUND__FOURTEEN_C2_FIELDS_BLOCKING__RELEASED_BENCHMARK_RUNTIME_ABSENT__ZERO_OF_SIX_PANEL_READY__PERFORMANCE_AND_SUPERIORITY_CANNOT_CHECK}.
The authoritative commit/tree/archive are pinned, but the tree has zero
\texttt{benchmark/} paths although the native runner requires
\texttt{benchmark/claw-eval}. No model, benchmark, protected case, scorer,
performance outcome or superiority result was executed or opened.

The distinct outcome-blind C3 DGM execution-binding successor is archived
under \idt{development/p5-dgm-execution-binding-v4-2026-08-23/}. Its exact
source/tree/uncompressed-Git-archive identity, source-rights manifest, closed
native parser, terminal rules, empty fallback policy, fail-closed runner,
resource registry, candidate-case and custody schemas, negative ledger and
134-check audit bind 6/21 C3 fields and preserve 15 blockers. Its manifest
contains 23 entries. The exact terminal is
\idt{P5_C3_V4_DGM_SOURCE_NATIVE_PARSER_EMPTY_FALLBACK_AND_OUTER_WALLCLOCK_BOUND__UNPINNED_DEPENDENCIES_NOT_MISLABELED_AS_LOCK__FIFTEEN_C3_FIELDS_BLOCKING__ZERO_OF_SIX_PANEL_READY__PERFORMANCE_AND_SUPERIORITY_CANNOT_CHECK}.
The dependency declarations have zero exact pins and no lockfile. The packet
opens no released outcome payload and executes no DGM, model, task, benchmark,
Docker container, protected case, score, performance outcome or superiority
contrast.

The distinct outcome-blind C4 ADIAS execution-binding successor is archived
under \idt{development/p5-adias-execution-binding-v4-2026-08-23/}. Its exact
source/tree/prefixed-Git-archive identity, 1,578-file source manifest,
source-rights record, closed native parser and terminal rules, empty fallback
policy, fail-closed runner, resource and dependency-declaration registries,
candidate-case/custody schemas, negative ledger and 35-check audit bind 6/21 C4
fields and preserve 15 blockers. Its manifest contains 23 entries. The exact
terminal is
\idt{P5_C4_V4_ADIAS_SOURCE_TREE_NATIVE_PARSER_EMPTY_FALLBACK_AND_WALLCLOCK_BOUND__FIFTEEN_C4_FIELDS_BLOCKING__DEPENDENCY_LOCK_TASK_RIGHTS_ISOLATION_AND_CUSTODY_UNBOUND__ZERO_OF_SIX_PANEL_READY__PERFORMANCE_AND_SUPERIORITY_CANNOT_CHECK}.
Four VCS declarations remain unpinned and task rights/isolation/custody remain
unbound. The packet executes no ADIAS model, domain, benchmark, task episode,
protected case, score, performance outcome or superiority contrast.

The distinct outcome-blind C5 Double Ratchet metric-only execution-binding
successor is archived under
\idt{development/p5-double-ratchet-execution-binding-v4-2026-08-23/}. Its exact
source/tree/prefixed-Git-archive, source-rights and 113-blob metadata manifests,
native parser/terminal rules, isolated write surface, empty fallback policy,
resource envelope, 46-entry uv lock, custody schemas, negative ledger and
351-check audit bind 9/21 C5 fields and preserve 12 blockers. Its manifest
contains 24 entries. The exact terminal is
\idt{P5_C5_V4_DOUBLE_RATCHET_SOURCE_PARSER_ISOLATION_FALLBACK_WALLCLOCK_COMPUTE_AND_DEPENDENCY_LOCK_BOUND__TWELVE_C5_FIELDS_BLOCKING__OFFICIAL_RUNNER_REGENERATES_SOLVER_OUTPUTS_AND_REPORTS_DEVELOPMENT_LOCKED_EACH_ROUND__ZERO_OF_SIX_PANEL_READY__PERFORMANCE_AND_SUPERIORITY_CANNOT_CHECK}.
The official loop regenerates hosted solver outputs and exposes its
development-locked metric every round; no P5-native task adapter or external
one-shot panel is bound. The packet executes no model, solver-output set,
benchmark, study case, protected score, performance outcome or superiority
contrast.

The distinct outcome-blind C6 ScienceClaw execution-binding successor is
archived under
\idt{development/p5-scienceclaw-execution-binding-v4-2026-08-23/}. Its exact
source/tree identity, source-rights record, strict dry-run-draft parser,
terminal rules, outer wallclock, case/custody schemas and negative ledger bind
5/21 C6 fields and preserve 16 blockers. No native selector or close-all
fallback switch exists, and its thirteen prior-result-prefix files were
hashed/censused without decoding. The packet executes no investigation, model,
tool, service, benchmark, prior-result payload or protected outcome. Fresh
source-level validator assurance is \textsc{CannotCheck} after cleanup because
the required \texttt{.source-audit} checkout is absent.

The panel-wide root-cause packet is archived under
\idt{development/p5-panelwide-root-closure-v5-2026-08-23/}. Its protocol,
cross-arm equivalence registry, result, recursive negative ledger,
fail-closed metadata gate, 74-check validator and 12-entry checksum manifest
recompute 42/126 bound and 84/126 blocking with zero of six ready. Fourteen
equivalence classes are partitioned into five root work packages, reducing the
coordination queue from 84 entries to five while creating zero arm-field
bindings. The packet also retains the unmatched wallclock envelopes and the
C4/C6 assurance defects. It executes no arm, model, benchmark or protected
scorer and opens no protected outcome; H1--H4, performance and superiority
remain \textsf{CANNOT\_CHECK}.

The first root-closure successor is archived under
\idt{development/p5-common-visible-case-rights-v6-2026-08-23/}. Its exact
Lang-1 source manifest, provenance, rights manifest, six-arm acceptance
receipt, recursive ledger, 741-check validator and 22-entry checksum manifest
bind 12 shared fields and recompute 54/126 bound, 72/126 blocking and zero of
six ready. The public source bodies are candidate-visible development content;
no arm, model, protected scorer or outcome is executed. Any future adapter that
adds content outside the frozen common core must reopen the corresponding
rights field.

The byte-level environment successor is archived under
\idt{development/p5-native-task-environment-fanout-v7-2026-08-23/}. It binds
C1's exact offline setup and effective configuration, records five distinct
arm-specific residual bundles, and recomputes 55/126 bound, 71/126 blocking
and zero of six ready. Its validator reports one environment bound, five
blocking, 67 referenced bytes verified and 25 packet files hashed. It executes
no arm, model, benchmark, protected scorer or outcome.

The C3 candidate-safe environment and its outcome-free initialization stop are
archived under
\idt{development/p5-c3-native-environment-v8-2026-08-23/} and
\idt{development/p5-c3-outcome-free-successor-v9-2026-08-23/}. V8 binds the
69-member seed, six unchanged Lang-1 members, eight control gates and the exact
native-initialization failure while opening none of the 1,595 excluded payloads.
V9 retains the preregistration, exact-byte/AST receipt and impossibility witness
without materializing an initializer. Both leave
\texttt{runtime.task\_environment} blocking; no DGM, model, benchmark, scorer or
outcome is executed, and the prohibited C4 validator is not rerun.

The three non-aggregable C2 lawful-successor packets are archived under
\idt{development/p5-c2-lawful-native-byte-successor-v11-2026-08-24/},
\idt{development/p5-c2-lawful-native-byte-successor-v12-2026-08-24/} and
\idt{development/p5-c2-lawful-native-byte-successor-v13-2026-08-24/}. V11
binds only \texttt{runtime.task\_environment} for a distinct 8/21 successor;
V12 binds only candidate-visible bytes and task-content rights for a distinct
9/21 successor; V13 binds only container/generated-artifact rights for a
distinct 8/21 successor. V13 contains the complete nine-file SPDX~2.3 SBOM,
rights map, licence bundle, generated-artifact authority, exported-layer
manifest, constrained runtime receipt, image archive and remote-environment
destruction receipt. The additive read-only integration audit is
\idt{development/p5-c2-v11-v13-manuscript-integration-audit-2026-08-24/}.
Its byte-identical paper-bound copy is
\idt{evidence/P5_C2_V11_V13_INTEGRATION_AUDIT.json}.
Each packet's validator and \idt{SHA256SUMS} pass without pytest or repository
CI. The three states are not unioned; released MOSS remains 7/21 bound and 14
blocking, 0/6 arms are ready, and H1--H4, performance and superiority remain
\textsf{CANNOT\_CHECK}.

The mapping from every result-bearing sentence in this manuscript to the archived artifact that supports it is \idt{evidence/CLAIM_LEDGER_V1.json}, with a human-readable view at \idt{evidence/CLAIM_LEDGER_V1.md}. The local hostile falsifier suite is \idt{evidence/FALSIFIER_V1.md}; the protected hidden-cause suite is \idt{evidence/hidden-cause-suite/PROTECTED_SUITE_V1.json}; and the nearest-work dispositions that Section~\ref{sec:related} discusses in prose are recorded per mechanism in \idt{evidence/TABLE_P5_1_NEAREST_WORK.json}.

\subsection*{Revision-level confirmatory archive}

The executed 96-case panel is bound by
\idt{protocol/SELF\_ORION\_V3\_REVISION\_LEVEL\_PROTOCOL\_V1.json}. Its
human-readable result is
\idt{evidence/revision-level-v3/CONFIRMATORY\_EXECUTION\_RESULT\_2026-08-24.md},
and the machine-readable receipt with per-policy metrics and independent
cross-checks is
\idt{research/self-orion-v3/confirmatory/CONFIRMATORY\_EXECUTION\_RECEIPT\_2026-08-24.json}.
The archive records \texttt{NO\_TERMINAL\_UNDER\_FROZEN\_RULES} and
\texttt{grants\_scientific\_authority=false}; it must not be represented as a
positive H1 or superiority result.

\subsection*{Diagnostic attribution archive}

The raw per-case records are \idt{evidence/glm-5.2-attribution/results.jsonl}, the checked sidecar report is \idt{evidence/glm-5.2-attribution/report.json}, and the reproduction receipt is \idt{evidence/glm-5.2-attribution/REPRODUCTION_RECEIPT.json}. The generated table data are under \idt{evidence/tables/} and their rendered forms under \idt{manuscript/tables/}. Architecture diagrams and pseudocode for the controller, and the historical technical companions this design inherits from, are retained under \idt{research/technical-companions/}.

\subsection*{P5-RD-02 public-development factorial and archival audit}

The three distinct Defects4J identities, their preregistrations, frozen-input
receipts, runners, raw phase streams and results are under
\idt{development/p5-rd02-defects4j-public-factorial-2026-08-23/}.  The bounded
scientific interpretation and immutable adverse-history map are
\idt{SCIENTIFIC_REPORT.md} and \idt{NEGATIVE_RESULT_LEDGER.md}; the rights,
authority and stream-provenance audit is \idt{AUDIT_RECEIPT.json}, with every
retained artifact bound by \idt{SHA256SUMS}.  The V3 result digest is
\idt{e743ebb77a31c4c4606f7eded8045a60c78fa2f251c13908051156cfd4a1fb6e}.
These artifacts support one local public known-fix replay and an archival audit
only.  They contain no protected/fresh task and grant no H1--H4 authority.

\section{Reproducibility}
\label{sec:reproducibility}

Regeneration of every archived result, without executing a live system, is

\begin{verbatim}
make paper05-results
\end{verbatim}

\noindent which runs \texttt{python -m orion.study.p5.tables}. The command recomputes 21/24 from the raw records, writes the confusion matrix and the residual-error ledger, and writes stubs for the seven tables that have no admissible rows. A rewrite of the archive that hides the three residual errors is refused with an error exit; a missing archive is reported as undetermined under its own distinct exit code rather than as a build failure, so that an absent measurement cannot be read as a failed build or as a negative result. Exact commands are in \idt{papers/orion-15-self-orion/REPRODUCE.md}. Continuous integration regenerates these artifacts on a clean checkout and requires every committed file to be byte-identical to what the command produces.

\subsection*{Compute and conservative-abstention cost}

The diagnostic archive records 14\,190 tokens and a mean latency of $\approx$13\,s per case, for attribution only. Cost and time to \emph{protected validated improvement} remain undetermined. No campaign-level abstention/cost trade-off is reported.

\subsection*{Licensing}

This revision does not add a repository-level license. Redistribution terms therefore remain undetermined from the repository state until they are declared explicitly.

\section{Ethics and safety}
\label{sec:ethics}

This work studies governed method evolution for research software. The frozen local falsifier and the diagnostic attribution suite use constructed hidden-cause cases. They do not involve human-subject records or personally identifiable information. No live provider campaign was executed in this revision.

The controller has no self-merge primitive, and its strongest internal terminal is a host-only promotion recommendation. Isolated self-certification ablations, if ever executed, are forbidden on the production authority path. Negative, null, blocked and undetermined outcomes are retained rather than dropped.
